\begin{abstract}
This paper studies how regional governments allocate a fixed industrial-policy capacity between a locally high-return activity and an activity that is complementary across jurisdictions. Cross-regional matches create additional value, but the region that undertakes the complementary activity captures only part of that value. We characterize the full Nash equilibrium set and compare it with a coordinated fixed-capacity benchmark. In the baseline continuous portfolio game, coordinated reallocation begins at $A^P=\Delta$, whereas differentiated decentralized equilibria first appear at $A^N=\Delta/(1-\alpha)$. Because $A^P<A^N$, there is a positive-measure interval in which complete priority duplication is the unique Nash equilibrium even though reallocating the same total policy capacity raises aggregate modeled surplus. A binary switching representation isolates incomplete local capture as the source of the threshold wedge. Restricted constant-returns and alternative-matching exercises preserve the threshold logic while showing that complete regional specialization is baseline-specific. The result identifies a composition-specific coordination problem in regional industrial policy: coordination can improve the allocation of priorities without increasing total policy capacity.
\end{abstract}
